\documentclass[10pt,a4paper]{amsart}
\usepackage{amsmath,amssymb,amsthm,graphicx}
% For commutative diagrams we use Xy-pic
\usepackage[all]{xy}

% We use 2cell for 2-commutative diagrams.
\xyoption{2cell}
% Theorem environments.
%
\theoremstyle{plain}
\newtheorem{theorem}[subsection]{Theorem}
\newtheorem{proposition}[subsection]{Proposition}
\newtheorem{lemma}[subsection]{Lemma}

\theoremstyle{definition}
\newtheorem{definition}[subsection]{Definition}
\newtheorem{example}[subsection]{Example}
\newtheorem{exercise}[subsection]{Exercise}
\newtheorem{situation}[subsection]{Situation}

\theoremstyle{remark}
\newtheorem{remark}[subsection]{Remark}
\newtheorem{remarks}[subsection]{Remarks}

\numberwithin{equation}{subsection}

% Macros
%
\def\lim{\mathop{\mathrm{lim}}\nolimits}
\def\colim{\mathop{\mathrm{colim}}\nolimits}
\def\Spec{\mathop{\mathrm{Spec}}}
\def\Hom{\mathop{\mathrm{Hom}}\nolimits}
\def\Ext{\mathop{\mathrm{Ext}}\nolimits}
\def\SheafHom{\mathop{\mathcal{H}\!\mathit{om}}\nolimits}
\def\SheafExt{\mathop{\mathcal{E}\!\mathit{xt}}\nolimits}
\def\Sch{\mathit{Sch}}
\def\Mor{\mathop{\mathrm{Mor}}\nolimits}
\def\Ob{\mathop{\mathrm{Ob}}\nolimits}
\def\Sh{\mathop{\mathit{Sh}}\nolimits}
\def\NL{\mathop{N\!L}\nolimits}
\def\CH{\mathop{\mathrm{CH}}\nolimits}
\def\proetale{{pro\text{-}\acute{e}tale}}
\def\etale{{\acute{e}tale}}
\def\QCoh{\mathit{QCoh}}
\def\Ker{\mathop{\mathrm{Ker}}}
\def\Im{\mathop{\mathrm{Im}}}
\def\Coker{\mathop{\mathrm{Coker}}}
\def\Coim{\mathop{\mathrm{Coim}}}

% Boxtimes
%
\DeclareMathSymbol{\boxtimes}{\mathbin}{AMSa}{"02}

%
% Macros for moduli stacks/spaces
%
\def\QCohstack{\mathcal{QC}\!\mathit{oh}}
\def\Cohstack{\mathcal{C}\!\mathit{oh}}
\def\Spacesstack{\mathcal{S}\!\mathit{paces}}
\def\Quotfunctor{\mathrm{Quot}}
\def\Hilbfunctor{\mathrm{Hilb}}
\def\Curvesstack{\mathcal{C}\!\mathit{urves}}
\def\Polarizedstack{\mathcal{P}\!\mathit{olarized}}
\def\Complexesstack{\mathcal{C}\!\mathit{omplexes}}
% \Pic is the operator that assigns to X its picard group, usage \Pic(X)
% \Picardstack_{X/B} denotes the Picard stack of X over B
% \Picardfunctor_{X/B} denotes the Picard functor of X over B
\def\Pic{\mathop{\mathrm{Pic}}\nolimits}
\def\Picardstack{\mathcal{P}\!\mathit{ic}}
\def\Picardfunctor{\mathrm{Pic}}
\def\Deformationcategory{\mathcal{D}\!\mathit{ef}}
\usepackage{cleveref}
\crefname{originallemma}{original lemma}{original lemmas}
\Crefname{originallemma}{Original lemma}{Original lemmas}
\makeatletter
\newlabel{algebra-lemma-flat-increases-depth}{{00LM}{source}{}{sourcepointer-1}{}}
\newlabel{algebra-lemma-flat-increases-depth@cref}{{[originallemma][0][]00LM}{[1][0][]source}{}{}{}}
\makeatother
\crefname{originallemma}{original lemma}{original lemmas}
\Crefname{originallemma}{Original lemma}{Original lemmas}
\makeatletter
\newlabel{algebra-lemma-grothendieck}{{00MF}{source}{}{sourcepointer-2}{}}
\newlabel{algebra-lemma-grothendieck@cref}{{[originallemma][0][]00MF}{[1][0][]source}{}{}{}}
\makeatother
\crefname{originallemma}{original lemma}{original lemmas}
\Crefname{originallemma}{Original lemma}{Original lemmas}
\makeatletter
\newlabel{lemma-cap-cp-bivariant}{{0B7H}{source}{}{sourcepointer-3}{}}
\newlabel{lemma-cap-cp-bivariant@cref}{{[originallemma][0][]0B7H}{[1][0][]source}{}{}{}}
\makeatother
\crefname{originallemma}{original lemma}{original lemmas}
\Crefname{originallemma}{Original lemma}{Original lemmas}
\makeatletter
\newlabel{lemma-chow-ring-projective-bundle}{{02TX}{source}{}{sourcepointer-4}{}}
\newlabel{lemma-chow-ring-projective-bundle@cref}{{[originallemma][0][]02TX}{[1][0][]source}{}{}{}}
\makeatother
\crefname{originallemma}{original lemma}{original lemmas}
\Crefname{originallemma}{Original lemma}{Original lemmas}
\makeatletter
\newlabel{lemma-geometric-cap}{{02SQ}{source}{}{sourcepointer-5}{}}
\newlabel{lemma-geometric-cap@cref}{{[originallemma][0][]02SQ}{[1][0][]source}{}{}{}}
\makeatother
\title{The top Chern class of a finite locally free module of rank r caps the fundamental class to the cycle of a regular section}
\author{The Stacks Project Authors}
\date{stacks-project:chow.tex:local-snapshot-20261012; GNU FDL 1.2 or later}
\begin{document}
\maketitle
\noindent\textit{Editorial note.} All delivered mathematical text is the author's own, copied from the pinned Stacks Project chapter file chow.tex; no statement or proof step was rewritten, reordered or removed.\par\medskip
\noindent\textit{Editorial note.} Preamble provenance: the pinned chapter file opens with \textbackslash{}input\{preamble\} and carries no preamble text of its own, so the project's single Stacks-wide preamble staged at sources/stacks-chow/preamble.tex is excerpted verbatim as two spans, lines 14-18 (Xy-pic setup) and lines 154-217 (the native theorem declarations and mathematical macros). Nothing was added, removed or redefined; the wrapper class and the amsmath/amssymb/amsthm/graphicx base packages are supplied by the assembler.\par\medskip
\noindent\textit{Editorial note.} Same-chapter pointers are located at their \textbackslash{}label lines inside the pinned chapter file; cross-chapter pointers are located in the project's own tag list staged in the snapshot (sources/stacks-chow/tags.txt), at the verbatim lines carrying their canonical tags.\par\medskip
\noindent\textit{Allowed standard prerequisites (editorial).}\begin{itemize}
\item The Chern class c\_r(E) is a bivariant class. Same chapter chow.tex, label lemma-cap-cp-bivariant, invoked at proof line 8765. Allows the base change X' -> X to be absorbed, reducing the statement to the case X = X'; the rest of the argument stays in the counted span.
\item The geometric case of the top Chern class for a regular section. Same chapter chow.tex, label lemma-geometric-cap, invoked at proof lines 8770 and 8820. Supplies the case r = 1 and is applied again to prove the auxiliary claim [Z(t)]\_\{n+r-2\} = c\_1(O\_P(1)) cap [P]\_\{n+r-1\}.
\item The projective space bundle formula and depth along a flat morphism. Same chapter chow.tex, label lemma-chow-ring-projective-bundle, invoked at proof line 8779; Algebra, labels algebra-lemma-flat-increases-depth (proof line 8785) and algebra-lemma-grothendieck (proof line 8831). Lets the equality be checked after pulling back along the projective bundle of E, keeps the regular-sequence hypothesis under flat pullback, and supplies the local nonzerodivisor criterion used in the claim.
\end{itemize}\medskip
\ifdefined\hypertarget\hypertarget{sourcepointer-1}{}\fi
\noindent\textit{Source reference.} Original lemma 00LM: sources/stacks-chow/tags.txt, source lines 792--792. Navigation pointer for \textbackslash{}ref\{algebra-lemma-flat-increases-depth\} used inside the counted span (line 8785); the argument that consumes it is retained in the span.\par\medskip
\ifdefined\hypertarget\hypertarget{sourcepointer-2}{}\fi
\noindent\textit{Source reference.} Original lemma 00MF: sources/stacks-chow/tags.txt, source lines 821--821. Navigation pointer for \textbackslash{}ref\{algebra-lemma-grothendieck\} used inside the counted span (line 8831); the argument that consumes it is retained in the span.\par\medskip
\ifdefined\hypertarget\hypertarget{sourcepointer-3}{}\fi
\noindent\textit{Source reference.} Original lemma 0B7H: sources/stacks-chow/pinned-original-source.tex, source lines 7533--7533. Navigation pointer for \textbackslash{}ref\{lemma-cap-cp-bivariant\} used inside the counted span (line 8765); the argument that consumes it is retained in the span.\par\medskip
\ifdefined\hypertarget\hypertarget{sourcepointer-4}{}\fi
\noindent\textit{Source reference.} Original lemma 02TX: sources/stacks-chow/pinned-original-source.tex, source lines 6932--6932. Navigation pointer for \textbackslash{}ref\{lemma-chow-ring-projective-bundle\} used inside the counted span (line 8779); the argument that consumes it is retained in the span.\par\medskip
\ifdefined\hypertarget\hypertarget{sourcepointer-5}{}\fi
\noindent\textit{Source reference.} Original lemma 02SQ: sources/stacks-chow/pinned-original-source.tex, source lines 4432--4432. Navigation pointer for \textbackslash{}ref\{lemma-geometric-cap\} used inside the counted span (line 8770); the argument that consumes it is retained in the span.\par\medskip
\begin{lemma}
\label{lemma-top-chern-class}
In the situation described just above assume $\dim_\delta(X') = n$,
that $f^*\mathcal{E}$ has constant rank $r$, that
$\dim_\delta(Z(s)) \leq n - r$, and that for every generic point
$\xi \in Z(s)$ with $\delta(\xi) = n - r$ the ideal of $Z(s)$
in $\mathcal{O}_{X', \xi}$ is generated by a regular sequence
of length $r$. Then
$$
c_r(\mathcal{E}) \cap [X']_n = [Z(s)]_{n - r}
$$
in $\CH_*(X')$.
\end{lemma}

\begin{proof}
Since $c_r(\mathcal{E})$ is a bivariant class
(Lemma \ref{lemma-cap-cp-bivariant})
we may assume $X = X'$ and we have to show that
$c_r(\mathcal{E}) \cap [X]_n = [Z(s)]_{n - r}$ in $\CH_{n - r}(X)$.
We will prove the lemma by induction on $r \geq 0$. (The case
$r = 0$ is trivial.) The case $r = 1$
is handled by Lemma \ref{lemma-geometric-cap}. Assume $r > 1$.

\medskip\noindent
Let $\pi : P \to X$ be the projective space bundle associated to
$\mathcal{E}$ and consider the short exact sequence
$$
0 \to \mathcal{E}' \to \pi^*\mathcal{E} \to \mathcal{O}_P(1) \to 0
$$
By the projective space bundle formula
(Lemma \ref{lemma-chow-ring-projective-bundle})
it suffices to prove the equality after pulling back by $\pi$.
Observe that $\pi^{-1}Z(s) = Z(\pi^*s)$ has $\delta$-dimension
$\leq n - 1$ and that the assumption on regular sequences at
generic points of $\delta$-dimension $n - 1$ holds by
flat pullback, see
Algebra, Lemma \ref{algebra-lemma-flat-increases-depth}.
Let $t \in \Gamma(P, \mathcal{O}_P(1))$ be the image of $\pi^*s$.
We claim
$$
[Z(t)]_{n + r - 2} = c_1(\mathcal{O}_P(1)) \cap [P]_{n + r - 1}
$$
Assuming the claim we finish the proof as follows.
The restriction $\pi^*s|_{Z(t)}$ maps to zero in
$\mathcal{O}_P(1)|_{Z(t)}$ hence comes from a unique
element $s' \in \Gamma(Z(t), \mathcal{E}'|_{Z(t)})$.
Note that $Z(s') = Z(\pi^*s)$ as closed subschemes of $P$.
If $\xi \in Z(s')$ is a generic point with $\delta(\xi) = n - 1$,
then the ideal of $Z(s')$ in $\mathcal{O}_{Z(t), \xi}$
can be generated by a regular sequence of length $r - 1$: it is generated by
$r - 1$ elements which are the images of $r - 1$ elements in
$\mathcal{O}_{P, \xi}$ which together with a generator of the
ideal of $Z(t)$ in $\mathcal{O}_{P, \xi}$ form a regular sequence
of length $r$ in $\mathcal{O}_{P, \xi}$. Hence we can apply the
induction hypothesis to $s'$ on $Z(t)$ to get
$c_{r - 1}(\mathcal{E}') \cap [Z(t)]_{n + r - 2} = [Z(s')]_{n - 1}$.
Combining all of the above we obtain
\begin{align*}
c_r(\pi^*\mathcal{E}) \cap [P]_{n + r - 1}
& =
c_{r - 1}(\mathcal{E}') \cap c_1(\mathcal{O}_P(1)) \cap [P]_{n + r - 1} \\
& =
c_{r - 1}(\mathcal{E}') \cap [Z(t)]_{n + r - 2} \\
& =
[Z(s')]_{n - 1} \\
& = [Z(\pi^*s)]_{n - 1}
\end{align*}
which is what we had to show.

\medskip\noindent
Proof of the claim. This will follow from an application of
the already used Lemma \ref{lemma-geometric-cap}.
We have $\pi^{-1}(Z(s)) = Z(\pi^*s) \subset Z(t)$.
On the other hand, for $x \in X$ if $P_x \subset Z(t)$, then
$t|_{P_x} = 0$ which implies that $s$ is zero in the fibre
$\mathcal{E} \otimes \kappa(x)$, which implies $x \in Z(s)$.
It follows that $\dim_\delta(Z(t)) \leq n + (r - 1) - 1$.
Finally, let $\xi \in Z(t)$ be a generic point with
$\delta(\xi) = n + r - 2$. If $\xi$ is not the generic point
of the fibre of $P \to X$ it is immediate that
a local equation of $Z(t)$ is a nonzerodivisor in $\mathcal{O}_{P, \xi}$
(because we can check this on the fibre by
Algebra, Lemma \ref{algebra-lemma-grothendieck}).
If $\xi$ is the generic point of a fibre, then $x = \pi(\xi) \in Z(s)$
and $\delta(x) = n + r - 2 - (r - 1) = n - 1$. This is a contradiction
with $\dim_\delta(Z(s)) \leq n - r$ because $r > 1$
so this case doesn't happen.
\end{proof}

\end{document}
